% Standalone machine note.  The Human Note is a separate, authoritative file.
\documentclass[11pt]{article}
\usepackage[T1]{fontenc}
\usepackage[utf8]{inputenc}
\usepackage{amsmath,amssymb,mathtools}
\usepackage{booktabs,graphicx}
\usepackage{xcolor}
\usepackage[margin=1in]{geometry}
\usepackage{microtype}
\usepackage{enumitem}
\usepackage[colorlinks=true,linkcolor=blue!55!black,urlcolor=blue!55!black]{hyperref}
\hypersetup{
  pdftitle={Physical X+psi denominator from the Human-Note plumbing sum},
  pdfauthor={Machine Note}
}
\allowdisplaybreaks
\setlist{nosep}
\title{Physical \(X+\psi\) denominator from the Human-Note plumbing sum}
\author{Machine Note}
\date{25 August 2026}

\begin{document}
\maketitle
\tableofcontents
\bigskip

Status: theta- and glasses-channel resummations derived and implemented.  The
noncompact loop-momentum Gaussian is left as a direct plumbing Gaussian and
is \textbf{not} replaced here by a canonical-frame spectral determinant.  This
note concerns the physical free boson and physical Majorana in the denominator
of \(Q_L\), not the auxiliary Majorana used in the double-Virasoro quotient.
The Human Note is not modified.

\section{Physical basis and the literal nonchiral sum}

Use the theta edge order

\[
 e=(0,1,\infty).
\]

For one chiral physical free superfield, use the ordered Fock basis

\[
 |p;\lambda,S\rangle
 =a_{-\lambda_1}\cdots a_{-\lambda_k}
   \psi_{-r_1}\cdots\psi_{-r_s}|p\rangle,
 \qquad
 r_1>\cdots>r_s>0,
\]

where \(\lambda\) is an integer partition and
\(S\subset\mathbb Z_{>0}-\tfrac12\).  Define

\[
 N(\lambda,S)=|\lambda|+\sum_{r\in S}r,
 \qquad
 \pi(\lambda,S)=|S|\pmod2,
 \qquad h(p)=\frac{p^2}{2}.
\]

Write \(U_e=(\lambda_e,S_e)\), \(V_e=(\mu_e,T_e)\) for the two
ends selected by the inverse Gram matrix, with tilded letters for the
antiholomorphic copy.  After dividing by the connected target-space volume,
the physical momentum measure is

\[
 [d\mathbf p]_{\Theta}
 =\prod_{e=0,1,\infty}dp_e\;
   \delta(p_0+p_1+p_\infty),
\]

up to the fixed convention-dependent powers of \(2\pi\).  The literal
Human-Note theta prescription becomes

\[
\boxed{
\begin{aligned}
Z_{X+\psi}^{\Theta,\mathrm{pl}}
={}&\int[d\mathbf p]_{\Theta}
\sum_{\mathbf U,\mathbf V,
      \widetilde{\mathbf U},\widetilde{\mathbf V}}
\prod_e
 q_e^{h(p_e)+N(U_e)}
 \bar q_e^{h(p_e)+N(\widetilde U_e)}
 \eta_e^{\pi(U_e)}
 \widetilde\eta_e^{\pi(\widetilde U_e)}
\\
&\quad\times
(-1)^{Q(\boldsymbol\pi(\mathbf U)
          +\boldsymbol\pi(\widetilde{\mathbf U}))}
\prod_e
 \left(\mathbb B_e^{-1}\right)^{
 U_e\widetilde U_e;V_e\widetilde V_e}
\\
&\quad\times
\mathcal R_{\Theta}(\mathbf U,\widetilde{\mathbf U};\mathbf p)
\mathcal R_{\Theta}(\mathbf V,\widetilde{\mathbf V};\mathbf p),
\end{aligned}}
\tag{1.1}
\]

where

\[
 Q(x_0,x_1,x_\infty)
 =x_0x_1+x_0x_\infty+x_1x_\infty\pmod2.
\tag{1.2}
\]

Here \(\mathbb B_e^{-1}\) is the inverse Gram matrix of the \textbf{full
nonchiral physical} \(X+\psi\) state.  Equation (1.1), rather than a theta
constant or a spectral determinant, is the definition of the denominator in
the theta plumbing Weyl frame.

The sign in (1.1) is not yet a product of a holomorphic and an
antiholomorphic sign.  The inverse tensor-product Gram matrix and the two
full pants correlators must be factorized before making such a statement.

\section{What the analytic free-field three-point function implies}

Normalize

\[
 X(Z_1)X(Z_2)\sim-\log Z_{12},\qquad
 Z_{ij}=z_i-z_j-\theta_i\theta_j,
\]

and let

\[
 \mathcal V_p(Z)=:e^{ipX(Z)}:,
 \qquad h(p)=\frac{p^2}{2}.
\]

The holomorphic primary three-point function is

\[
 \left\langle
 \mathcal V_{p_\infty}(Z_\infty)
 \mathcal V_{p_1}(Z_1)
 \mathcal V_{p_0}(Z_0)
 \right\rangle
 =\delta(p_0+p_1+p_\infty)
  \prod_{i<j}Z_{ij}^{p_ip_j}.
\tag{2.1}
\]

It contains the even superconformal three-point structure and no term
proportional to the odd invariant \(\Theta_{123}\).  Therefore

\[
 C^{(1)}_{p_0p_1p_\infty}=0.
\tag{2.2}
\]

This fact is what makes the free-theory reduction simpler than the
super-Liouville numerator.  It must be used only after the full Human-Note
sign algebra.  Here is that factorization explicitly.

Set

\[
 r_e=\pi(U_e)=\pi(V_e),
 \qquad
 \widetilde r_e=\pi(\widetilde U_e)=\pi(\widetilde V_e).
\]

The equalities follow because the inverse Gram matrix connects states at the
same level and hence with the same fermion parity.  The physical exponential
primaries are even.  Specializing the Human-Note tensor-product BPZ rule to
these free primaries gives

\[
 \left(\mathbb B_e^{-1}\right)^{
 U_e\widetilde U_e;V_e\widetilde V_e}
 =(-1)^{r_e\widetilde r_e}
  \left(B_e^{-1}\right)^{U_eV_e}
  \left(\widetilde B_e^{-1}\right)^{
    \widetilde U_e\widetilde V_e}.
\tag{2.2a}
\]

Each full pants correlator separately has a Koszul sign from moving all
antiholomorphic operators to the right.  The two pants correlators have the
same edge parities, so these two signs square to one.  With the normalized
free primary coefficient stripped off, their product is therefore

\[
\begin{aligned}
&\mathcal R_\Theta(\mathbf U,\widetilde{\mathbf U};\mathbf p)
 \mathcal R_\Theta(\mathbf V,\widetilde{\mathbf V};\mathbf p)
\\
&\qquad=
 \rho(\mathbf U;\mathbf p)\rho(\mathbf V;\mathbf p)
 \widetilde\rho(\widetilde{\mathbf U};\mathbf p)
 \widetilde\rho(\widetilde{\mathbf V};\mathbf p).
\end{aligned}
\tag{2.2b}
\]

Consequently the exponent of the complete sign in (1.1) is

\[
 \epsilon_{\mathrm{tot}}
 =Q(\mathbf r+\widetilde{\mathbf r})
  +\mathbf r\mathbin{\cdot}\widetilde{\mathbf r}.
\tag{2.2c}
\]

For three edges the elementary identity

\[
 Q(\mathbf r+\widetilde{\mathbf r})
 +\mathbf r\mathbin{\cdot}\widetilde{\mathbf r}
 =Q(\mathbf r)+Q(\widetilde{\mathbf r})
  +(r_0+r_1+r_\infty)
   (\widetilde r_0+\widetilde r_1+\widetilde r_\infty)
 \pmod2
\tag{2.2d}
\]

does all the remaining work.  The holomorphic and antiholomorphic pants
selection rules require

\[
 a=r_0+r_1+r_\infty
  =\widetilde r_0+\widetilde r_1+\widetilde r_\infty
 \pmod2.
\tag{2.2e}
\]

The last term of (2.2d) is therefore \(a^2=a\).  Thus the full nonchiral
sign becomes

\[
 (-1)^{a+Q(\mathbf r)+Q(\widetilde{\mathbf r})}.
\tag{2.3}
\]

Because (2.2) forces \(a=0\), the complete nonchiral sum is

\[
 Z_{X+\psi}^{\Theta,\mathrm{pl}}
 =\int[d\mathbf p]_{\Theta}
  \left|
  \mathcal F_{X+\psi,0}^{\Theta,\mathrm{pl}}
       (\mathbf p;\mathbf q,\boldsymbol\eta)
  \right|^2.
\tag{2.4}
\]

Thus holomorphic factorization is a result of (2.1)--(2.3), not an
assumption made in (1.1).  In particular, the descendant sign inside the
chiral block remains

\[
 (-1)^{Q(\pi_0,\pi_1,\pi_\infty)}.
\tag{2.5}
\]

\section{Analytic physical-Majorana pants tensor}

Let the positive integer \(m\) label the NS mode

\[
 \psi_{-(m-1/2)}|0\rangle.
\]

Use slot order \((\infty,1,0)\), and write a combined index as
\(\alpha=(s,m)\).  The analytic sphere contraction matrix is the
antisymmetric matrix \(K\) whose independent entries are

\[
\begin{aligned}
K_{(\infty,m),(1,n)}
 &=\begin{cases}
   \binom{m-1}{n-1},&m\ge n,\\[2mm]0,&m<n,
   \end{cases}\\
K_{(\infty,m),(0,n)}&=\delta_{mn},\\
K_{(1,m),(0,n)}
 &=(-1)^{m-1}\binom{m+n-2}{m-1}.
\end{aligned}
\tag{3.1}
\]

For three Fock states \(S_\infty,S_1,S_0\), let
\(I=S_\infty\sqcup S_1\sqcup S_0\) be the corresponding set of combined
indices.  In the generating-function order where modes increase on all three
slots, Wick's theorem gives

\[
 \widehat\rho_\psi(S_\infty,S_1,S_0)=\operatorname{Pf}K_I.
\tag{3.2}
\]

The physical BPZ bra reverses its fermion fields.  Consequently the physical
pants coefficient used in the literal state sum is

\[
 \rho_\psi(S_\infty,S_1,S_0)
 =(-1)^{|S_\infty|(|S_\infty|-1)/2}
  \widehat\rho_\psi(S_\infty,S_1,S_0).
\tag{3.3}
\]

Both coefficients vanish when the total number of fermions is odd.  In the
theta graph the two identical pants coefficients are multiplied, so the BPZ
reversal sign in (3.3) squares to one.  In the glasses graph it is retained
by the half-edge pairing ledger in Section 7.

\section{Fredholm resummation of the Human-Note descendant sign}

Introduce the plumbing half-variables

\[
 x_e=\eta_e\sqrt{q_e}.
\tag{4.1}
\]

For \(\boldsymbol\sigma=(\sigma_\infty,\sigma_1,\sigma_0)\in
\{\pm1\}^3\), define the diagonal mode-weight operator

\[
 (W_{\boldsymbol\sigma})_{(e,m),(e,m)}
 =\sigma_e x_e^{2m-1},
\tag{4.2}
\]

and the Fredholm determinant

\[
 D_{\boldsymbol\sigma}
 =\det_{\mathrm F}(1+K W_{\boldsymbol\sigma}).
\tag{4.3}
\]

The principal-minor identity and (3.2) give

\[
 D_{\boldsymbol\sigma}
 =\sum_{S_\infty,S_1,S_0}
   \widehat\rho_\psi(S_\infty,S_1,S_0)^2
   \prod_e\sigma_e^{|S_e|}
   x_e^{2N(S_e)}.
\tag{4.4}
\]

Only parity patterns \((000,110,101,011)\) occur.  On these four patterns,
the Human-Note sign (2.5) is respectively \((+,-,-,-)\).  Its exact
all-level projection is consequently

\[
\boxed{
 \mathcal F_{\psi}^{\Theta,\mathrm{pl}}
 =\frac12\left[
 -D_{(+,+,+)}
 +D_{(-,+,+)}
 +D_{(+,-,+)}
 +D_{(+,+,-)}
 \right].}
\tag{4.5}
\]

This is the desired physical-fermion resummation.  It uses the plumbing
coordinates, the physical lifts, and the analytic pants tensor only.  No
period matrix occurs.

At a mode cutoff \(M\), (4.5) requires four determinants of size \(3M\).
The cost is \(O(M^3)\), instead of enumerating an exponentially growing
collection of Fock states.  Numerically one should use

\[
 \det(1+KW)=\det(1+W^{1/2}KW^{1/2})
\tag{4.6}
\]

to balance the growing binomial entries of \(K\) against the decaying
plumbing weights.

\section{Charged boson and complete one-superfield resummation}

The bosonic oscillator pants sum exponentiates by Wick's theorem.  In the
same theta plumbing marking its all-level resummation is

\[
 P_X^{\Theta,\mathrm{pl}}(\mathbf q)
 =\prod_{[\gamma]\in\mathcal P_{\mathrm{prim}}}
   \prod_{n=1}^{\infty}(1-k_\gamma^n)^{-1},
\tag{5.1}
\]

where every multiplier \(k_\gamma\) is computed from the Schottky generators
constructed from the same \((0,1,\infty)\) plumbing maps.  Equation (5.1) is a
resummation of the Heisenberg pants state sum; it is not a canonical
Laplacian determinant.

The charged tensor is just as explicit.  In the bra convention

\[
 \langle\alpha_\infty|V_{\alpha_1}(1)|\alpha_0\rangle,
 \qquad \alpha_\infty=\alpha_1+\alpha_0,
\]

divide every current mode by the square root of its mode number.  In slot
order \((\infty,1,0)\), the symmetric normalized current-current kernel
\(K_X\) is obtained from

\[
\begin{aligned}
C_{\infty1}(m,n)&=
 m\binom{m-1}{n-1}\mathbf1_{m\ge n},\\
C_{\infty0}(m,n)&=m\delta_{mn},\\
C_{10}(m,n)&=(-1)^{m-1}n\binom{n+m-1}{m-1}
\end{aligned}
\tag{5.2}
\]

by \((K_X)_{(e,m),(f,n)}=C_{ef}(m,n)/\sqrt{mn}\).  The normalized
current-exponential contractions form the linear vector

\[
 \ell_{\infty,m}=\frac{\alpha_1}{\sqrt m},\qquad
 \ell_{1,m}=\frac{(-1)^{m-1}\alpha_0}{\sqrt m},\qquad
 \ell_{0,m}=-\frac{\alpha_1}{\sqrt m}.
\tag{5.3}
\]

Thus the analytic charged pants generating function is

\[
 \exp\!\left(\frac12t^TK_Xt+\ell^Tt\right).
\tag{5.4}
\]

This formula includes both pairwise current contractions and the singleton
contractions of a current with one of the exponential primaries.  Expanding
(5.4) gives the charged Fock-module tensor
\(\rho_{\alpha_\infty,\alpha_1,\alpha_0}(\lambda,\mu,\nu)\).

At mode cutoff \(M\), let

\[
 W_{(e,m),(e,m)}=q_e^m,\qquad D=W^{1/2},\qquad B=DK_XD.
\tag{5.5}
\]

The inverse-BPZ sewing convention gives the same numerical charge labels to
the two pants tensors.  Contracting the two Gaussian generating functions
therefore gives

\[
 P_{X,M}^{\Theta,\mathrm{pl}}
 =\det(1-B^2)^{-1/2},
\qquad
 E_M(\alpha_0,\alpha_1)
 =(D\ell)^T(1-B)^{-1}(D\ell).
\tag{5.6}
\]

The branch of the determinant is fixed by \(P_{X,M}\to1\) at the plumbing
cusp.  The complete fixed-charge chiral boson block is

\[
 \mathcal F_{X,M}^{\Theta,\mathrm{pl}}
 =P_{X,M}^{\Theta,\mathrm{pl}}
  \prod_{e=0,1,\infty}q_e^{\alpha_e^2/2}
  \exp E_M(\alpha_0,\alpha_1),
 \qquad
 \alpha_\infty=-(\alpha_0+\alpha_1)
\tag{5.7}
\]

when all three propagating charges are written as outgoing charges.  The
bra label in (5.3) is then \(-\alpha_\infty=\alpha_0+\alpha_1\).

Equations (5.5)--(5.7) are the requested direct resummation of the charged
free-boson state sum.  In particular, \(P_{X,M}\) converges to (5.1) without
using a period matrix.

The nonchiral charge dependence is a two-dimensional Gaussian.  Define its
quadratic form directly by

\[
 \left|
 \frac{\mathcal F_X^{\Theta,\mathrm{pl}}
       (\alpha_0,\alpha_1;\mathbf q)}
      {P_X^{\Theta,\mathrm{pl}}(\mathbf q)}
 \right|^2
 =\exp\!\left[-\pi
 (\alpha_0,\alpha_1)A_\Theta^{\mathrm{pl}}(\mathbf q)
 \binom{\alpha_0}{\alpha_1}\right].
\tag{5.8}
\]

The entries of \(A_\Theta^{\mathrm{pl}}\) are read directly from (5.7), or
equivalently from its values at charge vectors \((1,0),(0,1),(1,1)\).  With
completeness measure \(d\alpha_0d\alpha_1\),

\[
 \mathcal G_X^{\Theta,\mathrm{pl}}
 =\int_{\mathbb R^2}d^2\alpha\,
   e^{-\pi\alpha^TA_\Theta^{\mathrm{pl}}\alpha}
 =\det A_\Theta^{\mathrm{pl}}{}^{-1/2}.
\tag{5.9}
\]

No period matrix is used in (5.2)--(5.9).  A rescaling of the momentum
variable must rescale both \(A_\Theta^{\mathrm{pl}}\) and the completeness
measure.  This makes visible the otherwise hidden factor-of-two distinction
between a vertex charge with \(h=\alpha^2/2\) and a commonly used momentum
variable \(p=\sqrt2\alpha\) with \(h=p^2/4\).

Combining (4.5), (5.1), and (5.9) gives the period-matrix-free plumbing
formula for one physical free superfield:

\[
\boxed{
 Z_{X+\psi}^{\Theta,\mathrm{pl}}
 =\mathcal G_X^{\Theta,\mathrm{pl}}
  \left|P_X^{\Theta,\mathrm{pl}}\right|^2
  \left|\mathcal F_{\psi}^{\Theta,\mathrm{pl}}\right|^2.}
\tag{5.10}
\]

The \,\(\widehat c=9\) physical denominator is

\[
 Z_{\mathrm{free},\widehat c=9}^{\Theta,\mathrm{pl}}
 =\left(Z_{X+\psi}^{\Theta,\mathrm{pl}}\right)^9.
\tag{5.11}
\]

No auxiliary double-Virasoro fermion occurs in (1.1)--(5.11).

\section{Relation to the theta-constant shortcut}

If one separately proves, in precisely this plumbing trivialization and with
the Human-Note lifts, that

\[
 \left(\mathcal F_\psi^{\Theta,\mathrm{pl}}\right)^2
 =\vartheta[\delta](0|\Omega_{\mathrm{native}})
  P_X^{\Theta,\mathrm{pl}},
\tag{6.1}
\]

then (5.10) can be algebraically rewritten in the earlier theta-constant
form.  Equation (6.1) is not assumed in the construction above.  It is now a
cross-check between two resummations, so a possible determinant-line or
marking mistake cannot silently define the physical denominator.

\section{Glasses-channel resummation}

For the glasses edge order \((L,R,B)\), direct chiral sewing gives

\[
\begin{aligned}
\mathcal F_\psi^{\mathrm{Gl},\mathrm{pl}}
={}&\sum_{S_L,S_R,S_B}
 (-1)^{\pi_B(\pi_L+\pi_R)}
 \prod_{e=L,R,B}x_e^{2N(S_e)}\\
&\qquad\times
 \rho_\psi(S_L,S_B,S_L)
 \rho_\psi(S_R,S_B,S_R).
\end{aligned}
\tag{7.1}
\]

Each nonzero pants tensor requires an even total number of fermions.  Since
the self-loop state occurs twice, this forces \(\pi_B=0\).  Thus the explicit
Human-Note glasses sign in (7.1) is one on the support of the free-Majorana
tensor.  It has not been discarded before this selection rule is applied.

Let

\[
 K_{\psi,\Gamma}=K\oplus K,
 \qquad
 w_{e,m}=x_e^{2m-1},
\tag{7.2}
\]

and put \(\sqrt{w_{e,m}}\) on each of the two half-edges incident on edge
\(e\); their diagonal matrix is \(D_\psi\).  Let \(J_\psi\) be the
antisymmetric half-edge pairing matrix.  In the half-edge order
\((\infty_L,1_L,0_L,\infty_R,1_R,0_R)\), its nonzero edge orientations are

\[
 J_{\psi,L}=-1,
 \qquad J_{\psi,R}=-1,
 \qquad J_{\psi,B}=i.
\tag{7.3}
\]

The two minus signs are the self-loop bends.  The bridge phase is the
coefficient-basis correction: contracting two Grassmann generating functions
gives \((-1)^{s(s-1)/2}\), while the surviving bridge occupations have even
\(s\).  Therefore the factor \(i^s=(-1)^{s/2}\) in (7.3) cancels it.  The
choice \(-i\) gives the same result because odd bridge occupation vanishes.
The exact cutoff resummation is

\[
\boxed{
 \mathcal F_{\psi,M}^{\mathrm{Gl},\mathrm{pl}}
 =\det\!\left(
  1+J_\psi D_\psi K_{\psi,\Gamma}D_\psi
 \right)^{1/2}.}
\tag{7.4}
\]

The square-root branch is the one approaching one at the plumbing cusp.
Equation (7.4) was checked against the literal state sum (7.1), not inferred
from a theta constant.

For the boson, set \(K_{X,\Gamma}=K_X\oplus K_X\), put
\(\sqrt{q_e^m}\) on each half-edge in \(D_X\), and let \(J_X\) be the
symmetric pairing with entry one across all three edges.  Then

\[
 P_{X,M}^{\mathrm{Gl},\mathrm{pl}}
 =\det\!\left(1-J_XD_XK_{X,\Gamma}D_X\right)^{-1/2}.
\tag{7.5}
\]

For self-loop charges \((\alpha_L,\alpha_R)\), the concatenated source is
\(\ell_\Gamma=(\ell(\alpha_L,0),\ell(\alpha_R,0))\).  With
\(j=D_X\ell_\Gamma\), the charged exponent and primary propagation are

\[
 E_M^{\mathrm{Gl}}
 =\frac12j^T
 \left(1-J_XD_XK_{X,\Gamma}D_X\right)^{-1}J_Xj,
 \qquad
 q_L^{\alpha_L^2/2}q_R^{\alpha_R^2/2}.
\tag{7.6}
\]

The bridge primary is neutral.  Reading the real quadratic form
\(A_{\mathrm{Gl}}^{\mathrm{pl}}\) from the absolute square of (7.6) gives

\[
 \mathcal G_X^{\mathrm{Gl},\mathrm{pl}}
 =\det A_{\mathrm{Gl}}^{\mathrm{pl}}{}^{-1/2},
\tag{7.7}
\]

in the displayed \(d\alpha_Ld\alpha_R\) measure.  Hence

\[
\boxed{
 Z_{X+\psi}^{\mathrm{Gl},\mathrm{pl}}
 =\mathcal G_X^{\mathrm{Gl},\mathrm{pl}}
  \left|P_X^{\mathrm{Gl},\mathrm{pl}}\right|^2
  \left|\mathcal F_\psi^{\mathrm{Gl},\mathrm{pl}}\right|^2.}
\tag{7.8}
\]

Both (5.10) and (7.8) are raw plumbing-frame quantities.  Their numerical
values should not be equated without the channel-change/Weyl factors that
belong to the full comparison.

\section{Unpacking every factor in the physical denominator}

The compact notation in (5.10) and (7.8) contains three logically different
objects.  For one real physical superfield they are

\[
 \underbrace{\mathcal G_X^{\mathrm{pl}}}_{
   \text{two physical loop-charge integrals}}
 \quad\times\quad
 \underbrace{|P_X^{\mathrm{pl}}|^2}_{
   \text{nonzero bosonic oscillators}}
 \quad\times\quad
 \underbrace{|\mathcal F_\psi^{\mathrm{pl}}|^2}_{
   \text{physical NS-Majorana oscillators}}.
\tag{8.1}
\]

The last factor is the physical partner \(\psi\) of \(X\).  It is not the
auxiliary Majorana introduced in the double-Virasoro branching identity.

\subsection{Theta graph}

At mode cutoff \(M\), define

\[
 \Delta_{\psi,\Theta,M}
 =-D_{(+,+,+),M}+D_{(-,+,+),M}
   +D_{(+,-,+),M}+D_{(+,+,-),M},
\tag{8.2}
\]

where each \(D_{\boldsymbol\sigma,M}\) is the \(3M\)-dimensional determinant
in (4.3), evaluated in its balanced form (4.6).  Then

\[
 |\mathcal F_{\psi,M}^{\Theta,\mathrm{pl}}|^2
 =\frac14|\Delta_{\psi,\Theta,M}|^2.
\tag{8.3}
\]

For the physical boson, let

\[
 B_M=D_MK_{X,M}D_M,
 \qquad
 D_M^2=\operatorname{diag}_{(e,m)}(q_e^m).
\tag{8.4}
\]

The nonzero-mode factor is therefore

\[
 |P_{X,M}^{\Theta,\mathrm{pl}}|^2
 =\left|\det(1-B_M^2)\right|^{-1}.
\tag{8.5}
\]

It remains to unpack the zero-mode Gaussian.  Put
\(u=(\alpha_0,\alpha_1)^T\), and introduce the charge-incidence matrix

\[
 R_\Theta=
 \begin{pmatrix}
  1&0\\0&1\\-1&-1
 \end{pmatrix},
 \qquad
 R_\Theta u=(\alpha_0,\alpha_1,\alpha_\infty)^T.
\tag{8.6}
\]

Let \(\Lambda_{\Theta,M}\) be the \(3M\)-by-2 matrix defined by
\(\ell=\Lambda_{\Theta,M}u\).  Its rows in slot order
\((\infty,1,0)\) are

\[
 (\Lambda_{\Theta,M})_{(\infty,m)}
  =\frac1{\sqrt m}(0,1),\qquad
 (\Lambda_{\Theta,M})_{(1,m)}
  =\frac1{\sqrt m}((-1)^{m-1},0),\qquad
 (\Lambda_{\Theta,M})_{(0,m)}
  =\frac1{\sqrt m}(0,-1).
\tag{8.7}
\]

Writing \(L_{\Theta,M}=D_M\Lambda_{\Theta,M}\), the charged oscillator
exponent in (5.6) is

\[
 E_M(u)=u^TH_{\Theta,M}u,
 \qquad
 H_{\Theta,M}
 =L_{\Theta,M}^T(1-B_M)^{-1}L_{\Theta,M}.
\tag{8.8}
\]

Here and below only the symmetric part of \(H\) contributes to \(u^THu\).
Taking the absolute square of the primary propagation and (8.8) gives the
loop matrix explicitly:

\[
\boxed{
 A_{\Theta,M}^{\mathrm{pl}}
 =-\frac1\pi\left[
 R_\Theta^T
 \operatorname{diag}(\log|q_0|,\log|q_1|,\log|q_\infty|)
 R_\Theta
 +2\operatorname{Re}\operatorname{Sym}H_{\Theta,M}
 \right].}
\tag{8.9}
\]

Thus the charge integral is not an independent canonical-frame input:

\[
 \mathcal G_{X,M}^{\Theta,\mathrm{pl}}
 =\det A_{\Theta,M}^{\mathrm{pl}}{}^{-1/2}.
\tag{8.10}
\]

Substituting (8.3), (8.5), and (8.10) completely removes the shorthand from
the one-superfield answer:

\[
\boxed{
 Z_{X+\psi,M}^{\Theta,\mathrm{pl}}
 =\frac{|\Delta_{\psi,\Theta,M}|^2}
 {4\,\sqrt{\det A_{\Theta,M}^{\mathrm{pl}}}\,
  |\det(1-B_M^2)|}.}
\tag{8.11}
\]

For nine independent physical superfields,

\[
\boxed{
 Z_{\mathrm{free},\widehat c=9,M}^{\Theta,\mathrm{pl}}
 =2^{-18}
  \bigl(\det A_{\Theta,M}^{\mathrm{pl}}\bigr)^{-9/2}
  |\det(1-B_M^2)|^{-9}
  |\Delta_{\psi,\Theta,M}|^{18}.}
\tag{8.12}
\]

All determinants in (8.11)--(8.12) are constructed from the theta plumbing
coordinates, lifts, and analytic pants kernels.

\subsection{Glasses graph}

Define the two cutoff matrices

\[
 \mathbb M_{\psi,\mathrm{Gl},M}
 =1+J_\psi D_\psi K_{\psi,\Gamma,M}D_\psi,
 \qquad
 \mathbb M_{X,\mathrm{Gl},M}
 =1-J_XD_XK_{X,\Gamma,M}D_X.
\tag{8.13}
\]

Equations (7.4) and (7.5) then give

\[
 |\mathcal F_{\psi,M}^{\mathrm{Gl},\mathrm{pl}}|^2
 =|\det\mathbb M_{\psi,\mathrm{Gl},M}|,
 \qquad
 |P_{X,M}^{\mathrm{Gl},\mathrm{pl}}|^2
 =|\det\mathbb M_{X,\mathrm{Gl},M}|^{-1}.
\tag{8.14}
\]

Put \(u=(\alpha_L,\alpha_R)^T\).  Let
\(\Lambda_{\mathrm{Gl},M}\) be the \(6M\)-by-2 matrix determined by

\[
 \ell_\Gamma
 =\bigl(\ell(\alpha_L,0),\ell(\alpha_R,0)\bigr)
 =\Lambda_{\mathrm{Gl},M}u,
 \qquad L_{\mathrm{Gl},M}=D_X\Lambda_{\mathrm{Gl},M}.
\tag{8.15}
\]

The matrix representing the exponent (7.6) is

\[
 H_{\mathrm{Gl},M}
 =\frac12L_{\mathrm{Gl},M}^T
  \mathbb M_{X,\mathrm{Gl},M}^{-1}J_X
  L_{\mathrm{Gl},M}.
\tag{8.16}
\]

The bridge primary is neutral, so the glasses loop matrix is

\[
\boxed{
 A_{\mathrm{Gl},M}^{\mathrm{pl}}
 =-\frac1\pi\left[
 \operatorname{diag}(\log|q_L|,\log|q_R|)
 +2\operatorname{Re}\operatorname{Sym}H_{\mathrm{Gl},M}
 \right].}
\tag{8.17}
\]

Consequently the fully unpacked glasses expressions are

\[
\boxed{
 Z_{X+\psi,M}^{\mathrm{Gl},\mathrm{pl}}
 =\frac{|\det\mathbb M_{\psi,\mathrm{Gl},M}|}
 {\sqrt{\det A_{\mathrm{Gl},M}^{\mathrm{pl}}}\,
  |\det\mathbb M_{X,\mathrm{Gl},M}|},}
\tag{8.18}
\]

and

\[
\boxed{
 Z_{\mathrm{free},\widehat c=9,M}^{\mathrm{Gl},\mathrm{pl}}
 =\bigl(\det A_{\mathrm{Gl},M}^{\mathrm{pl}}\bigr)^{-9/2}
  |\det\mathbb M_{\psi,\mathrm{Gl},M}|^9
  |\det\mathbb M_{X,\mathrm{Gl},M}|^{-9}.}
\tag{8.19}
\]

\subsection{Normalization ledger}

Equations (8.11)--(8.19) use the following precise normalization:

\begin{enumerate}
\item \(XX\sim-\log z\) and \(h(\alpha)=\alpha^2/2\).
\item The two independent genus-two charges are integrated with
\(d^2\alpha\), after dividing out the connected target-space volume.
\item The normalized free-field pants coefficient is one.
\item The plumbing propagator is \(q^{L_0}\); the common cylinder Casimir
factor \(q^{-c/24}\) is stripped.
\item No channel-change Weyl factor, determinant-line trivialization factor, or
auxiliary double-Virasoro fermion is included.
\end{enumerate}

If the completeness convention is \(d\alpha/(2\pi)\) for every real charge,
there are two independent charges for each of nine superfields, and the
right-hand sides of (8.12) and (8.19) acquire the overall factor

\[
 (2\pi)^{-18}.
\tag{8.20}
\]

If instead the propagator convention includes the cylinder vacuum energy,
one real \(X+\psi\) superfield has \(c=3/2\).  Its nonchiral multiplier is

\[
 \prod_{e=1}^3|q_e|^{-c/12}
 =\prod_{e=1}^3|q_e|^{-1/8},
\tag{8.21}
\]

and the nine-superfield multiplier is

\[
 \prod_{e=1}^3|q_e|^{-9/8}.
\tag{8.22}
\]

These optional constants must be applied consistently to both channels;
they are not hidden inside \(A\), \(P_X\), or \(\mathcal F_\psi\).
If the Human-Note definition of \(L_e\) already contains the Casimir shift,
(8.21)--(8.22) must not be inserted a second time.

At the \texttt{\detokenize{o0243-periodmatched}} point, the numerical multiplication of the
three normalized factors is

\begin{center}
\small
\resizebox{\textwidth}{!}{%
\begin{tabular}{lrrrrr}
\toprule
\textbf{channel} & \textbf{\(\mathcal G_X\)} & \textbf{\(|P_X|^2\)} & \textbf{\(|\mathcal F_\psi|^2\)} & \textbf{\(Z_{X+\psi}\)} & \textbf{\(Z_{X+\psi}^9\)} \\
\midrule
theta & 0.215006349267368 & 1.00000166309736 & 1.00258782998132 & 0.215563107646015 & \(1.00500893632052\,10^{-6}\) \\
glasses & 1.16308106016403 & 1.33846154450516 & 2.69778130216318 & 4.19974210080833 & \(4.06446695711553\,10^5\) \\
\bottomrule
\end{tabular}%
}
\end{center}

These are raw plumbing-frame values with the normalization ledger above.
Their inequality is not a failure of the resummation: the channel-change and
Weyl factors have deliberately not been applied.

\section{Recomputed \texorpdfstring{\(Q_L\)}{Q L} with the plumbing denominator}

The existing Human-Note-sign super-Liouville numerators at
\texttt{\detokenize{o0243-periodmatched}} can be reused: changing the
independent physical free-theory denominator does not change a conformal
block, structure constant, or momentum-quadrature node.  Replacing the old
period-matrix shortcut by (8.11) and (8.18), at the common cutoff \(M=24\),
gives

\begin{center}
\small
\resizebox{\textwidth}{!}{%
\begin{tabular}{lrrrr}
\toprule
\textbf{channel} & \textbf{direct \(Z_{X+\psi}^{\rm pl}\)} &
\textbf{old shortcut} & \textbf{direct/old} &
\textbf{\(Q_L\) multiplier} \\
\midrule
theta & 0.215563107646015 & 0.431128572561718 &
0.499997266163927 & 512.025195722 \\
glasses & 4.19974210080833 & 8.39948420161306 &
0.500000000000214 & 511.999999998 \\
\bottomrule
\end{tabular}%
}
\end{center}

The factor is fixed by the displayed charge convention.  If
\(q=\exp(2\pi i\Omega)\), then for \(h(\alpha)=\alpha^2/2\),

\[
 \int_{\mathbb R^g}d^g\alpha\,
 e^{-2\pi\alpha^T\operatorname{Im}\Omega\,\alpha}
 =2^{-g/2}\det(\operatorname{Im}\Omega)^{-1/2}.
\tag{9.1}
\]

At \(g=2\), this is the factor \(1/2\) that was absent from the old displayed
shortcut.  Raising one physical superfield to the ninth power therefore
multiplies the absolute \(Q_L\) by approximately \(2^9=512\).  The small
theta displacement from exactly 512 is the previously observed difference
between the finite plumbing resummation and the shortcut, not a failure of
the cutoff convergence.

The recomputed values are

\begin{center}
\small
\begin{tabular}{rrrr}
\toprule
\(R\) & \(N\) & \(Q_L^\Theta/Q_L^{\rm Gl}\) &
\(Q_L^\Theta/Q_L^{\rm Gl}-1\) \\
\midrule
0 & 4 & 1.08934489273 & 0.08934489273 \\
3 & 4 & 1.08996087620 & 0.08996087620 \\
4 & 4 & 1.09091076293 & 0.09091076293 \\
4 & 6 & 1.04478367268 & 0.04478367268 \\
\bottomrule
\end{tabular}
\end{center}

For the best local setting \((R,N)=(4,6)\),

\[
 Q_L^\Theta=7.23725849554\times10^{-12},\qquad
 Q_L^{\rm Gl}=6.92704019478\times10^{-12}.
\tag{9.2}
\]

A subsequent full local rerun of all four \((R,N)\) settings with the direct
denominator wired into the evaluator reproduced every numerator and every
recombined \(Q_L\) exactly as stored in
\texttt{\detokenize{Data Set/ns_genus2_theta_glasses_hatc9_human_note_plumbing_free_2026-08-25.json}}.

Thus the new physical-Majorana and charged-boson plumbing computation
changes the absolute normalization substantially but does not remove the
channel discrepancy: the current ratio differs from one by \(4.478\%\).
The free denominator itself is stable well beyond \(10^{-6}\): from
\(M=20\) to \(M=24\), the relative changes are below \(3\times10^{-15}\) in
theta and \(2\times10^{-13}\) in glasses.  The remaining visible movement is in
the super-Liouville numerator, especially its momentum quadrature.

\section{Audit and removal of the apparent four-percent discrepancy}

The completed corrected-lift calculation at \((R,N)=(24,10)\) first appeared
to give

\[
 \frac{Q_L^\Theta}{Q_L^{\rm Gl}}=1.04056522832
 \tag{10.1}
\]

after replacing the old denominator by the direct plumbing denominator.
This is not a free-denominator cutoff effect.  The changes from \(M=20\) to
\(M=24\) are below \(2\times10^{-13}\), and replacing the direct denominator
by the old period-matrix shortcut moves (10.1) by only
\(4.9\times10^{-5}\) relatively.  The available cutoff axes likewise show
that the \(R=24\) recursion is stable and that the remaining \(N=10\)
quadrature movement is at the few-times-\(10^{-4}\) scale, not four percent.

The origin is instead a coefficient-convention mismatch at the interface
between the real BRY super-Liouville data and the graded factorization in the
Human Note.  This is fixed by comparing definitions, not by fitting (10.1).
For positive continuum momenta, the implementation
\texttt{ns\_tilde\_structure\_constant} returns BRY's real coefficient of the
locally normalized top component

\[
 W_{\rm BRY}=G_{-1/2}\widetilde G_{-1/2}V,
 \qquad
 \langle W_{\rm BRY},W_{\rm BRY}\rangle=+D(2h)^2,
\]

where $D$ is the matched primary two-point normalization.  The Human Note
instead sews the ordered tensor state

\[
 w_{\rm HN}=(G_{-1/2}V)\otimes
 (\widetilde G_{-1/2}\widetilde V).
\]

It declares left- and right-moving odd operators to anticommute and uses the
graded tensor-product BPZ pairing.  Therefore its definition gives directly

\[
 \langle w_{\rm HN},w_{\rm HN}\rangle_{\rm HN}=-D(2h)^2.
\]

After matching the primary normalization and the chiral three-form
normalization, the boundary dictionary is consequently

\[
 C_{\rm HN}^{(0)}=C_{\rm BRY},\qquad
 C_{\rm HN}^{(1)}=\sigma i\,\widetilde C_{\rm BRY},
 \qquad \sigma\in\{+1,-1\}.
\tag{10.2}
\]

Thus the BPZ comparison fixes the squared map, but not its square-root branch.
The implementation chooses \(\sigma=+1\); \(\sigma=-1\) gives the same
genus-two result because every term contains two pants coefficients.  The odd
coefficient product supplies

\[
 (C_{\rm HN}^{(1)})^2=-\widetilde C_{\rm BRY}^{2}
\]

in theta, and the same minus sign in
\(C^{(1)}_{LBL}C^{(1)}_{RBR}\) in glasses.  This sign is \emph{not} a change
to the Human Note.  Its explicit descendant/Koszul sign \((-1)^a\) remains
exactly as written; (10.2) only converts the imported BRY coefficient into
that convention before the Human-Note formula is applied.

This scope distinction is essential.  BRY-native Type~0B sphere correlators,
including the $G,H,J$ formulas and PCO amplitudes, retain the real
\(\widetilde C_{\rm BRY}\) with no additional $i$.  Equation~(10.2) is
applied exactly once only when BRY data enter a Human-Note graded sewing sum.
It is also unrelated to both the physical free Majorana in the denominator
of \(Q_L\) and the auxiliary Majorana in the double-Virasoro construction.
The removal of the four-percent discrepancy corroborates this dictionary;
it is not the argument that determines it.  Fixing \(\sigma\) in an amplitude
with an odd number of odd pants would require one explicitly ordered
component matrix element.

The \((24,10)\) shard decomposition makes the diagnosis numerical and
unambiguous:

\begin{center}
\small
\begin{tabular}{lrr}
\toprule
channel & even contribution & signed odd contribution before (10.2) \\
\midrule
theta & \(7.27241679849612\times10^{-18}\) & \(-2.77708574019467\times10^{-20}\) \\
glasses & \(2.88444939646816\times10^{-6}\) & \(-6.87808460834071\times10^{-8}\) \\
\bottomrule
\end{tabular}
\end{center}

Inserting (10.2) reverses the displayed odd contributions while retaining
the Human-Note sewing sign.  With the direct plumbing denominator this gives

\[
 \begin{aligned}
 Z_L^\Theta&=7.30018765589807\times10^{-18},&
 Q_L^\Theta&=7.26380372559184\times10^{-12},\\
 Z_L^{\rm Gl}&=2.95323024255156\times10^{-6},&
 Q_L^{\rm Gl}&=7.26597183274289\times10^{-12},
 \end{aligned}
\]

and therefore

\[
 \boxed{\frac{Q_L^\Theta}{Q_L^{\rm Gl}}
 =0.999701608099,\qquad
 \frac{Q_L^\Theta}{Q_L^{\rm Gl}}-1
 =-2.98391901\times10^{-4}.}
\tag{10.3}
\]

Thus the convention repair removes \(99.26\%\) of the previous residual.
It also improves monotonically in the inexpensive local quadrature check:

\begin{center}
\small
\begin{tabular}{rrr}
\toprule
\(R\) & \(N\) & corrected \(Q_L^\Theta/Q_L^{\rm Gl}-1\) \\
\midrule
4 & 4 & \(+4.84740320\times10^{-2}\) \\
4 & 6 & \(+3.83177997\times10^{-3}\) \\
4 & 8 & \(+1.40461919\times10^{-3}\) \\
24 & 10 & \(-2.98391901\times10^{-4}\) \\
\bottomrule
\end{tabular}
\end{center}

The remaining \(2.98\times10^{-4}\) is consistent with the unresolved
momentum-quadrature tail at \(N=10\); it is no longer evidence for a
four-percent plumbing, free-theory, or recursion failure.  A fresh common
\(R=24\), \(N=12\) or higher run is still required before claiming
\(10^{-6}\) stabilization.

\section{Executable implementation and checks}

The theta and glasses fermion determinants, charged-boson Gaussians, and
complete one- and nine-superfield expressions are implemented in

\begin{verbatim}
Code/genus_2/physical_free_plumbing_resummation.py
\end{verbatim}

The focused test compares (4.5) and (7.4) with direct physical-Majorana
Fock/Pfaffian sums for three lift assignments.  It also compares (5.6) and
(7.5) with the independent primitive-word boson products:

\begin{verbatim}
Code/genus_2/test_physical_free_plumbing_resummation.py
\end{verbatim}

All eleven focused checks pass.  The implementation reports separately the
fermion determinant, Heisenberg determinant, charged-boson Schur complement,
two-by-two loop matrix, one complete superfield, and its ninth power.

The fast production entry point

\begin{verbatim}
physical_superfield_plumbing_partition
\end{verbatim}

uses the bosonic Gaussian vacuum determinant directly.  Primitive-word
enumeration is retained only as an independent test oracle.

At the \texttt{\detokenize{o0243-periodmatched}} plumbing point, using only the listed plumbing
coordinates and Human-Note lifts, consecutive-cutoff stabilization below
\(10^{-6}\) occurs at

\begin{center}
\small
\resizebox{\textwidth}{!}{%
\begin{tabular}{lrr}
\toprule
\textbf{channel} & \textbf{common determinant cutoff} & \textbf{local time} \\
\midrule
theta & \(M=8\) & about 2 ms \\
glasses & \(M=10\) & about 10 ms \\
\bottomrule
\end{tabular}%
}
\end{center}

The converged raw ingredients in the \(h=\alpha^2/2\) charge convention are

\paragraph{Theta.}
\[
\begin{aligned}
\mathcal F_\psi&=1.00129307876105-2.01628153784\,10^{-5}i,\\
P_X&=1.00000083154833-3.00494124702\,10^{-8}i,\\
\mathcal G_X&=0.215006349267368,\\
Z_{X+\psi}&=0.215563107646015.
\end{aligned}
\]

\paragraph{Glasses.}
\[
\begin{aligned}
\mathcal F_\psi&=1.64249239061865+2.21185456733\,10^{-4}i,\\
P_X&=1.15691889842888-4.54930513585\,10^{-4}i,\\
\mathcal G_X&=1.16308106016403,\\
Z_{X+\psi}&=4.19974210080833.
\end{aligned}
\]

The timings are single-process local timings after import and are meant only
to show the computational order.  The production determinant computation
scales as \(O(M^3)\); primitive-word enumeration is no longer part of the
production denominator.
\end{document}
